\documentclass[conference]{../../../../Conference-LaTeX-template_10-17-19/IEEEtran}
\usepackage{cite}
\usepackage{amsmath,amssymb,amsfonts}
\usepackage{array}
\usepackage{booktabs}
\usepackage{tabularx}
\usepackage{url}
\newcolumntype{Y}{>{\raggedright\arraybackslash}X}
\renewcommand{\arraystretch}{1.12}

\title{Can High-Power Lasers Become Sunlight-Like?\\A Coherence-Resolved Research Proposal for Laser-Plasma Interaction Control}
\author{\IEEEauthorblockN{Anonymous Research Proposal}}

\begin{document}
\maketitle

\begin{abstract}
The possibility of a high-power laser whose statistical coherence resembles selected properties of sunlight has a precise scientific motivation: broadband spectrum, stochastic phase, and time-varying polarization can shorten resonance memory in laser-plasma interactions (LPI). Such control may reduce the growth of selected parametric instabilities that scatter driver energy or generate hot electrons, both of which can compromise inertial confinement fusion (ICF) coupling. However, the phrase ``sunlight-like laser'' is often broader than the available evidence. A published model and scoped particle-in-cell comparison support a mechanism for mitigation of stimulated Raman scattering (SRS) in a homogeneous plasma, while polarization-smoothing studies support fast-time speckle averaging. Neither result establishes a reactor-scale driver, a universal instability threshold, or a completed ignition improvement. This paper develops the \emph{Coherence-Resolved Instability Suppression Protocol for ICF} (CRISP-ICF), a source-bounded, design-only research plan. CRISP-ICF represents a candidate driver through spectral support, phase correlation, polarization statistics, temporal-speckle statistics, spatial-intensity statistics, and inter-beam independence. It compares coherent, spectral-smoothing, polarization-smoothing, vector-light, sunlight-like, and adversarial overlapping-beam cards through a staged analytic, wave/kinetic, and coupling-validation ladder. The proposal's central contribution is a multi-objective decision ledger: a candidate may advance only when a declared LPI-risk proxy improves under matched controls without an unacceptable penalty in drive uniformity, useful energy delivery, or cross-beam sensitivity. The report provides falsifiers, outcome branches, an error/evidence ledger, and human-review gates. It reports no new simulation or experimental result and does not provide facility-operation instructions.
\end{abstract}

\begin{IEEEkeywords}
laser-plasma interaction, inertial confinement fusion, broadband laser, polarization smoothing, stimulated Raman scattering, stimulated Brillouin scattering, coherence engineering, high-energy-density physics
\end{IEEEkeywords}

\section{Introduction}
Inertial confinement fusion requires a driver, a target, and an implosion pathway to work together under exceptionally demanding conditions. In that setting, the drive field is not just a source of energy. Its temporal, spectral, spatial, and polarization structure can influence how energy propagates through plasma, how nonuniformity is imprinted, and whether parasitic wave couplings draw energy into undesirable channels. The U.S. National Ignition Facility (NIF) reports repeated fusion-ignition results and an approved project to increase laser energy, demonstrating that high-energy-density laser fusion remains an active frontier rather than a hypothetical application \cite{nif_2026}. That official context does not establish that a particular coherence-engineered driver is ready for target-scale use, but it makes the underlying LPI question consequential.

The motivating question is therefore worth making exact: \emph{can humans make intense lasers with incoherence comparable to sunlight?} The answer has two parts. First, high-power laser systems can in principle be designed to have selected stochastic optical properties analogous to sunlight, including broad spectral support, randomized phase, and time-dependent polarization. They are not literal artificial suns, nor do they need to reproduce the full solar spectrum, angular extent, temporal statistics, brightness, or natural propagation history. Second, whether those properties are useful for ICF depends on a specified plasma regime and on system-level tradeoffs. A favorable reduction in one instability proxy is not sufficient if it produces an adverse spatial-intensity tail, a cross-beam coupling penalty, unacceptable delivery loss, or a driver architecture that cannot be realized and characterized at the required scale.

Published work provides a concrete mechanism rather than a metaphor. Ma \emph{et al.} formulate a sunlight-like field with continuous broad spectrum, random phase spectrum, and random polarization \cite{ma2021}. In their stated homogeneous-plasma numerical comparison, adjacent temporal speckles have limited duration and shifted central frequencies; the associated loss of resonance memory, coupled with random polarization, weakens the three-wave coupling that underlies selected parametric instabilities. Their model reported a higher SRS threshold for a roughly 1\% relative-bandwidth sunlight-like case than for a monochromatic comparator. The result is important, but its correct interpretation is narrow: it motivates a falsifiable driver-control hypothesis, not a facility-independent claim of ignition improvement.

This paper simulates the repository's four-agent scientific workflow:
\begin{center}
Survey $\rightarrow$ Idea $\rightarrow$ ExperimentDesign $\rightarrow$ Author.
\end{center}

The Survey Agent freezes evidence and gaps. The Idea Agent selects a design direction that is more discriminating than a bandwidth-only comparison. The ExperimentDesign Agent produces a computational-digital protocol with no automatic execution. The Author Agent organizes only those inputs into an IEEE-style research proposal. The report directly answers the topic while retaining the distinctions required for a credible research plan: what has been modeled, what has been verified from sources, what is proposed, and what still needs specialist review.

\section{Evidence Survey: What Is Established and What Is Not}
\subsection{Coherence engineering has several independent control dimensions}
The baseline for many LPI discussions is a highly coherent driver with a narrow spectral description, a defined polarization, and a focal pattern that can seed intense speckles. Smoothing by spectral dispersion (SSD) adds temporal variation to the speckle pattern by using bandwidth; polarization smoothing (PS) combines distinct orthogonally polarized patterns that add incoherently. Rothenberg's analysis of PS for ICF explains why the instantaneous incoherent addition can be especially relevant when an instability responds rapidly, and why PS can complement SSD rather than simply duplicate it \cite{rothenberg2000}. This establishes a useful starting point: a driver should not be described by bandwidth alone.

The sunlight-like formulation extends that idea by treating spectrum, phase, and polarization as a joint stochastic object. A broad spectrum can reduce persistent phase matching, random phase can prevent a temporally stable field realization, and random polarization can alter the coupling geometry. In an idealized notation, a two-component field may be represented as
\begin{equation}
\mathbf{E}(t)=\Re\left\{\sum_{j\in\{y,z\}}\mathbf{e}_j\int A_j(\omega)\exp\left[i\omega t+i\phi_j(\omega)\right]d\omega\right\},
\label{eq:field}
\end{equation}
where $A_j(\omega)$ is a component spectrum, $\phi_j(\omega)$ is a phase realization, and $\mathbf{e}_j$ is a polarization basis vector. Equation \eqref{eq:field} is a representation template, not a prescription for a physical driver. To be useful, it must be accompanied by normalization, component correlations, spatial propagation, and statistics evaluated at the intended interaction plane.

Polarization should also not be reduced to a scalar label. Vector-light work provides a separate, scoped mechanism in which spatially varying polarization can affect a side-scattering geometry \cite{jia2023}. It expands the candidate space but also reinforces the need for a comparison discipline: a polarization pattern that changes one scattering path cannot be assumed to improve every Raman, Brillouin, two-plasmon, or cross-beam process in every density profile.

\subsection{LPI mitigation is inherently a tradeoff problem}
Parametric instabilities are coupled processes. In an ICF-relevant plasma, SRS, SBS, and cross-beam energy transfer (CBET) can coexist with evolving density gradients, flows, multiple beams, and changing spatial statistics. A model that reduces a local growth proxy can still be inadequate if it changes overlap statistics or the distribution of high-intensity speckles. Hao \emph{et al.} provide a useful adversarial example: in their scoped simulations of overlapping smoothed beams, coupling changed the intensity distribution and could enhance SBS through high-intensity speckles \cite{hao2023}. This counterexample is central to the proposed method. It means that a single-beam, single-instability success condition cannot certify multi-beam usefulness.

Cross-beam modeling raises a second fidelity issue. Myatt \emph{et al.} motivate a wave-based treatment of CBET that retains polarization, plasma inhomogeneity, beam geometry, and realistic beam properties rather than relying only on a steady-state ray description \cite{myatt2017}. The present report treats that paper as a source-bounded model-fidelity motivation; its publisher-page review remains pending. The methodological implication is nevertheless clear: a candidate driver cannot cross from a reduced homogeneous screen to an ICF compatibility claim without an explicit propagation and coupling bridge.

\subsection{System relevance does not remove claim boundaries}
ICF has a long history of coupled driver, target, hydrodynamics, and diagnostic challenges \cite{betti2016,atzeni2004}. The contribution of coherence engineering must therefore be stated at the level the evidence supports. A source may justify a mechanism, a modeled threshold, or a smoothing rationale. It does not justify a claim that a sunlight-like high-power laser has already delivered better target gain, solved implosion symmetry, enabled repeatable power-plant operation, or replaced the need for target physics and driver engineering. The distinction matters because it preserves a path for ambitious research without turning a plausible hypothesis into a premature conclusion.

\section{Research Gaps and Idea Selection}
The Survey stage identifies four accepted gaps. First, there is no generally shared design contract that connects spectral width, phase correlation, polarization statistics, temporal-speckle lifetime, and spatial-intensity statistics to a declared LPI quantity of interest. Second, suppression-oriented comparisons do not automatically establish preservation of usable drive, uniformity, or coupling. Third, reduced or homogeneous simulations lack a validated transfer contract for multi-beam propagation and target-scale conditions. Fourth, driver-engineering feasibility is not proven by a mechanism simulation.

The Idea Agent considered three routes. A bandwidth-only threshold map is easy to explain but collapses polarization and overlap into untracked nuisance variables. A polarization-only dashboard preserves a fast-time smoothing route but does not characterize spectral resonance memory. The selected primary idea, \emph{Coherence-Resolved Instability Suppression Protocol for ICF} (CRISP-ICF), treats the full coherence structure as a controlled experimental-design object. Its objective is not to nominate a universal winner. It is to decide whether a candidate has enough matched, scoped, and uncertainty-aware evidence to justify a higher-fidelity specialist benchmark.

Figure \ref{fig:crisp} summarizes the proposed decision flow. The editable source artifact is retained as \texttt{figures/crisp\_icf\_flow.drawio}. A candidate passes from a driver card through a representation gate, a multi-fidelity LPI stack, and a tradeoff ledger. A status can be favorable only in the restricted sense of \texttt{COHERENCE\_BENEFIT\_CANDIDATE}; it is not a measured suppression result or a statement about ignition.

\begin{figure*}[!t]
\centering
\fbox{\begin{minipage}{0.95\textwidth}
\centering\footnotesize
\textbf{CRISP-ICF decision flow}\\[5pt]
\begin{tabular}{c@{\quad$\longrightarrow$\quad}c@{\quad$\longrightarrow$\quad}c@{\quad$\longrightarrow$\quad}c}
\parbox{0.16\textwidth}{\centering\textbf{Driver cards}\\D0 coherent\\D1 spectral\\D2 polarization\\D3 sunlight-like\\D4 vector\\D5 overlap stress} &
\parbox{0.16\textwidth}{\centering\textbf{Representation gate}\\spectrum\\energy normalization\\phase statistics\\Stokes statistics\\focal statistics} &
\parbox{0.18\textwidth}{\centering\textbf{Physics ladder}\\analytic resonance screen\\wave or kinetic comparison\\coupling stress test} &
\parbox{0.17\textwidth}{\centering\textbf{Tradeoff ledger}\\LPI risk proxy\\delivery and uniformity\\cross-beam sensitivity\\uncertainty and resources}\\[14pt]
\multicolumn{4}{c}{\textbf{Status:} baseline, candidate benefit, tradeoff, coupling-sensitive, inconclusive, or invalid.}
\end{tabular}
\end{minipage}}
\caption{CRISP-ICF connects a coherence representation to a conditional evidence status. It is a research-design figure, not a measured beamline result.}
\label{fig:crisp}
\end{figure*}

\section{ExperimentDesign: A Coherence-Resolved Computational Protocol}
\subsection{Quantity of interest and comparison contract}
The ExperimentDesign Agent selects the computational-digital template. The unit of analysis is one \emph{regime card $\times$ driver card $\times$ matched control $\times$ quantity-of-interest card}. A quantity of interest $Q$ may be a regime-specific backscatter proxy, an SRS or SBS risk surrogate, a hot-electron-risk surrogate, a drive-uniformity measure, a useful-energy-delivery proxy, or a cross-beam sensitivity. The protocol does not preselect a favorable metric and does not treat a driver as successful because one uncontrolled scalar changes in a desired direction.

Each card must declare its physical scope, reference route, units, representation assumptions, and uncertainty fields. The driver-coherence descriptor is
\begin{equation}
\mathbf{c}=\left(\frac{\Delta\omega}{\omega_0},S(\omega),\tau_{\phi},\mathcal{P}_{\mathrm{Stokes}},\tau_{\mathrm{speckle}},\mathcal{I}_{\mathrm{spatial}},\mathcal{C}_{\mathrm{beam}}\right),
\label{eq:descriptor}
\end{equation}
where $\Delta\omega/\omega_0$ is relative spectral support, $S(\omega)$ is spectral shape, $\tau_{\phi}$ is a phase-correlation descriptor, $\mathcal{P}_{\mathrm{Stokes}}$ specifies polarization statistics, $\tau_{\mathrm{speckle}}$ is a temporal-speckle descriptor, $\mathcal{I}_{\mathrm{spatial}}$ records spatial-intensity statistics, and $\mathcal{C}_{\mathrm{beam}}$ records inter-beam dependence. This vector is intentionally not reduced to a universal ``incoherence score.'' Different components can influence different LPI channels, and the appropriate representation depends on the selected regime.

The central comparison rule is strict. A sunlight-like candidate `D3` must be compared to a coherent baseline `D0` with compatible energy normalization, central wavelength convention, focal-envelope definition, physical scope, and numerical-resolution policy. It must also be compared to intermediate controls `D1` and `D2` where possible. This prevents a compound intervention from receiving credit for an untracked change in intensity, geometry, spectral envelope, or solver behavior.

\subsection{Staged model cards}
Table \ref{tab:cards} defines the cards. `D5` is deliberately adversarial: it represents credible overlap and cross-beam conditions rather than a favorable single-beam idealization. `R0` through `R3` define increasing model fidelity. Passing one tier authorizes consideration at the next tier; it does not validate the higher tier by implication.

\begin{table*}[!t]
\caption{CRISP-ICF model cards and non-negotiable validation obligations. Every value would be a planned input or derived analysis field; this proposal supplies no observed values.}
\label{tab:cards}
\centering
\small
\begin{tabularx}{\textwidth}{lYYY}
\toprule
Card & Purpose & Required fields & Failure response \\
\midrule
D0 coherent baseline & Establish matched reference & carrier, envelope, polarization, focal statistics, energy normalization & \texttt{MODEL\_INVALID} if matching is absent \\
D1 spectral control & Separate bandwidth/phase effects & spectrum, bandwidth, phase process, smoothing time & \texttt{REPRESENTATION\_INVALID} \\
D2 polarization control & Separate incoherent polarization addition & component balance, mutual coherence, focal statistics & \texttt{TRADEOFF\_DOMINATED} or invalid \\
D3 sunlight-like candidate & Test coupled spectrum-phase-polarization hypothesis & independent phase processes, Stokes statistics, temporal speckles, spatial statistics & \texttt{INCONCLUSIVE} if descriptors are incomplete \\
D4 vector-light candidate & Probe geometry-dependent polarization route & polarization map, propagation model, scattering geometry & scope-limited comparison only \\
D5 overlap stress control & Expose cross-beam and high-speckle risks & overlap geometry, beam independence, gradient, coupling model & \texttt{COUPLING\_SENSITIVE} \\
R0--R3 fidelity ladder & Prevent unqualified transfer & analytic scope, wave/kinetic model, transport compatibility, uncertainty & withhold target-scale conclusion \\
\bottomrule
\end{tabularx}
\end{table*}

The lowest layer, `R0`, verifies the field representation. The generated or modeled spectrum, polarization process, energy normalization, and spatial statistics must match the declaration in Eq. \eqref{eq:descriptor}. A broadband field at injection is not automatically a broadband or random-polarization field at the interaction plane. The next layer, `R1`, is an analytic or linear resonance-memory screen. It can reject a candidate that has no plausible decorrelation timescale for the declared mode. It cannot quantify target gain or certify a facility outcome.

`R2` is a tightly scoped wave or kinetic comparison that uses identical control fields except for the coherence element being tested. The intended comparison includes `D0`, `D1`, `D2`, and `D3`, with `D4` used when a side-scattering geometry is relevant. An appropriate numerical study would document solver approximation, resolution behavior, boundary conditions, initialization, noise control, and reference model. The proposal does not run that study. `R3` is a specialist-reviewed transport and drive-compatibility gate that incorporates propagation, beam overlap, and any necessary target-physics bridge. It is explicitly separate from the lower-fidelity mechanism result.

\subsection{Error, evidence, and resource ledger}
For a declared quantity of interest, CRISP-ICF records a conservative claim ledger,
\begin{align}
\epsilon_Q^{\mathrm{claim}} &\lesssim \epsilon_{\mathrm{model}}+\epsilon_{\mathrm{input}}+\epsilon_{\mathrm{repr}}+\epsilon_{\mathrm{algorithm}} \nonumber\\
&+\epsilon_{\mathrm{sampling}}+\epsilon_{\mathrm{coupling}}+\epsilon_{\mathrm{data}}.
\label{eq:errorledger}
\end{align}
where the terms denote model-form or regime discrepancy, uncertain physical inputs, field representation mismatch, algorithmic approximation, stochastic or ensemble sampling, inter-beam/coupling transfer, and reference-data limitations. Equation \eqref{eq:errorledger} is an accounting structure rather than an assumption that errors are independent or exactly additive. A missing term blocks a strong claim; it is not silently set to zero.

The multi-objective output is represented by
\begin{align}
\mathbf{y}=\bigl(&Q_{\mathrm{LPI}},Q_{\mathrm{delivery}},Q_{\mathrm{uniformity}},Q_{\mathrm{coupling}}, \nonumber\\
&\epsilon_Q^{\mathrm{claim}},R_{\mathrm{compute}}\bigr).
\label{eq:outcomevector}
\end{align}
where $Q_{\mathrm{LPI}}$ is a declared instability-risk proxy, $Q_{\mathrm{delivery}}$ is a useful-energy-delivery proxy, $Q_{\mathrm{uniformity}}$ is a drive-uniformity proxy, $Q_{\mathrm{coupling}}$ is a cross-beam sensitivity proxy, and $R_{\mathrm{compute}}$ is a resource interval. The protocol does not set arbitrary weights to produce a single winner. It uses Pareto-style reporting: a candidate that improves one coordinate while degrading another must report the tradeoff, not hide it behind an aggregate score.

\subsection{Falsifiers and reproducibility obligations}
CRISP-ICF is designed to be refutable. The core hypothesis fails for a declared regime if the candidate advantage disappears after matching spectral envelope, energy, focal statistics, and physical scope; if the status reverses under a credible overlap card; if representation uncertainty is comparable to the claimed difference; or if an intermediate control explains the change without the full sunlight-like structure. A positive result from `R1` or `R2` also fails as a target-scale claim if the `R3` bridge is missing.

The future implementation would require a provenance record for every driver realization, random-seed policy, field-spectrum verification, Stokes-statistics verification, solver version, mesh or particle-resolution study, boundary conditions, normalization convention, and analysis script. It would report both stable and switching regions in the selected parameter space. The report's purpose is to specify those obligations before a result is generated, when they can still prevent an attractive but non-comparable conclusion.

\section{Expected Outcome Branches and Conditional Conclusions}
Table \ref{tab:branches} limits what the Author stage may conclude after a future specialist-reviewed execution. These are branches of a predeclared analysis plan, not results. In particular, \texttt{COHERENCE\_BENEFIT\_CANDIDATE} authorizes a higher-fidelity benchmark and a driver-engineering feasibility review; it does not mean that an intense sunlight-like driver has been built, that absorption improved in a facility, or that ignition became controllable.

\begin{table}[!t]
\caption{Conditional-status framework for CRISP-ICF. Each label is an authorized interpretation only after the stated checks are complete.}
\label{tab:branches}
\centering
\footnotesize
\begin{tabularx}{\columnwidth}{lY}
\toprule
Label & Authorized conditional interpretation \\
\midrule
\texttt{REPRESENTATION\_INVALID} & The modeled driver does not meet its declared spectrum, phase, polarization, spatial-statistics, or normalization contract. \\
\texttt{BASELINE\_ONLY} & Existing coherent or conventional smoothing controls are sufficient for the stated scope; no added candidate case is supported. \\
\texttt{COHERENCE\_BENEFIT\_CANDIDATE} & Matched, scoped evidence justifies a higher-fidelity specialist benchmark; no observed performance claim follows. \\
\texttt{TRADEOFF\_DOMINATED} & A favorable LPI proxy is outweighed by a declared delivery, uniformity, or resource penalty. \\
\texttt{COUPLING\_SENSITIVE} & The conclusion changes under plausible overlap or cross-beam assumptions and cannot be generalized. \\
\texttt{INCONCLUSIVE} & Fidelity, sensitivity, or uncertainty coverage does not support a direction. \\
\texttt{MODEL\_INVALID} & Units, regime scope, reference route, or observation/projection status fails validation. \\
\bottomrule
\end{tabularx}
\end{table}

The conditional branches permit a direct scientific answer. Humans can plausibly synthesize high-intensity driver fields with selected low-coherence properties, and published mechanism studies justify investigating them. Whether those fields should be considered for an ICF-relevant driver is conditional: the field must remain characterized after propagation, show a robust matched-control advantage on the selected LPI proxy, avoid an unacceptable tradeoff, and survive a credible multi-beam/coupling test. This is more ambitious than a generic caveat because it identifies exact evidence needed to advance the claim.

\section{From LPI Control to ICF-System Relevance}
\subsection{A bridge, not an extrapolation}
The pathway from a coherence-engineered laser field to ICF relevance has several interfaces. First, a source architecture must create and preserve the declared spectrum, phase, and polarization statistics at useful energy and pulse format. Second, optical transport must preserve or quantify changes in the field at the interaction plane. Third, plasma propagation must be modeled in the appropriate density-gradient and multi-beam regime. Fourth, any LPI change must be interpreted alongside drive uniformity, energy delivery, and target compatibility. A good result at one interface is not automatically a good result at the next.

This proposal makes that separation constructive. It gives driver engineers a target set of measurable statistics rather than an imprecise metaphor. It gives LPI theorists matched controls and an adversarial overlap case. It gives target and systems researchers a rule that reduced-model evidence can motivate, but cannot substitute for, a propagation/transport bridge. It also gives reviewers a negative outcome that is informative: an \texttt{INCONCLUSIVE}, \texttt{COUPLING\_SENSITIVE}, or \texttt{TRADEOFF\_DOMINATED} label can prevent resources from being committed to a driver concept whose apparent advantage is an artifact of a favorable comparison.

\subsection{Why bandwidth remains necessary but insufficient}
Broadening the spectrum can be valuable because it changes phase-matching and smoothing timescales. Ma \emph{et al.} specifically emphasize that bandwidth requirements can be challenging for conventional approaches and propose a compounded temporal-speckle and polarization mechanism \cite{ma2021}. That insight motivates the `D3` card. Yet a scalar bandwidth cannot describe whether the spectrum is flat, Gaussian, structured, independently generated across polarization components, or preserved during transport. It also cannot describe whether temporal decorrelation introduces a problematic intensity structure or whether overlapping beams share correlations.

CRISP-ICF therefore treats bandwidth as one coordinate in Eq. \eqref{eq:descriptor}. A future driver study may find that a modest spectral support combined with appropriate phase and polarization statistics has a better tradeoff than a wider but poorly controlled spectrum. It may also find the reverse. The proposed research design does not embed either answer; it makes the comparison testable and requires source- and model-bounded language when the answer is reported.

\subsection{Potential technology implications}
If a candidate repeatedly passes the conditional gates, its next research questions are architectural rather than rhetorical: Can a high-energy amplifier chain preserve the intended stochastic statistics? Can optical components tolerate the relevant spectral and polarization content? Can diagnostics recover the needed coherence descriptors at the interaction plane? Can the candidate integrate with beam smoothing, frequency conversion, pulse shaping, and target requirements without introducing a new dominant failure mode? Those questions require qualified driver, optical, plasma, and facility teams. They are not answered by the present proposal, and they should not be converted into unreviewed design instructions.

\section{Pre-Registered Analysis and Decision Audit}
The CRISP-ICF protocol has value only if its decision rule is declared before a favorable field realization or a convenient solver setting is selected. A future study should therefore register its quantity of interest, physical regime, driver-card definitions, matched-control rule, numerical fidelity tier, uncertainty representation, and escalation criteria before running the comparison. This is not an administrative addition. It distinguishes a coherence effect from post hoc selection of an attractive spectrum, polarization trace, output time window, or plasma subregion.

The required audit matrix is shown in Table \ref{tab:audit}. The matrix creates a clear separation between an engineering hypothesis and a conclusion. For example, a `D3` candidate that has a favorable analytic resonance-memory proxy but fails the field-representation gate is not a positive result waiting for better language; it is an invalid representation. Similarly, a kinetic comparison that favors `D3` only when beam overlap is excluded is a coupling-sensitive result, not a general multi-beam proposal. Each row identifies what must be stored, what it prevents, and how the author is permitted to write the conclusion.

\begin{table*}[!t]
\caption{Pre-registered CRISP-ICF decision audit. The table specifies future evidence obligations; it reports no computed or measured values.}
\label{tab:audit}
\centering
\small
\begin{tabularx}{\textwidth}{lYYY}
\toprule
Audit item & Required record & Risk prevented & Authorized response if incomplete \\
\midrule
Driver identity & Spectrum, normalization, phase/polarization process, spatial and temporal statistics at the interaction plane & A nominally sunlight-like source differs from its modeled field or from the matched baseline & \texttt{REPRESENTATION\_INVALID} \\
Scope identity & Density/temperature/gradient representation, geometry, duration convention, and LPI quantity of interest & A local mechanism result is retold as a universal ICF claim & Restrict claim to the declared scope \\
Numerical verification & Solver version, resolution policy, boundary conditions, convergence or uncertainty evidence, seed policy & Numerical artifact is confused with an incoherence benefit & \texttt{INCONCLUSIVE} \\
Control separation & D0, D1, D2, D3, and applicable D4 cards share normalization and physical scope & Bandwidth, polarization, or intensity changes are conflated & \texttt{BASELINE\_ONLY} or invalid \\
Adversarial coupling & D5 overlap/cross-beam card, correlation assumptions, and sensitivity switches & Single-beam smoothing narrative hides an overlap penalty & \texttt{COUPLING\_SENSITIVE} \\
Transfer gate & Explicit model-fidelity tier, missing-physics list, and specialist review record & Reduced-model output is treated as target-scale validation & Withhold transport or ignition inference \\
\bottomrule
\end{tabularx}
\end{table*}

The audit also makes negative outcomes useful. A \texttt{BASELINE\_ONLY} result can establish that a conventional smoothing route is adequate in the chosen regime. A \texttt{TRADEOFF\_DOMINATED} result can show that an otherwise attractive coherence signature is not worth a delivery or uniformity penalty. An \texttt{INCONCLUSIVE} result can reveal which measurement, model, or computational resource is actually limiting the next experiment or simulation. These are scientifically substantive outcomes because they refine the boundary between a mechanism demonstration and a driver-relevant design.

Finally, the protocol requires a source-status field for each comparison input and interpretation. Statements based on a publisher-verified source, a browser-verified facility page, cross-validated metadata, a future resource projection, and a locally generated result must not be merged into a single undifferentiated evidence category. The field is essential for an AI-assisted workflow: it prevents fluent prose from silently upgrading a source's evidentiary status.

\section{Risks, Limitations, and Human Review Requirements}
The proposal has several nontrivial limitations. The field representation in Eq. \eqref{eq:field} is a compact stochastic abstraction; it does not capture every amplifier, optical, transport, or nonlinear propagation effect. The analytic screen cannot replace kinetic treatment where kinetic effects are material. A particle-in-cell or wave calculation has its own numerical noise, boundary, resolution, and reduced-physics risks. Cross-beam and target-scale inference introduce further model discrepancy. The ledger in Eq. \eqref{eq:errorledger} makes these sources visible but does not eliminate them.

The strongest evidence source in this report is a publisher-verified mechanism study with a stated homogeneous-plasma simulation scope \cite{ma2021}. The NIF source is browser-verified official context \cite{nif_2026}. Rothenberg, Jia, and Hao were discovered and cross-validated through OpenAlex and AnySearch; their final publisher-page bibliographic review remains required before external dissemination \cite{rothenberg2000,jia2023,hao2023}. Myatt and the general ICF background sources are used only within the limited roles stated in the evidence appendix. This source discipline matters because a review paper, a simulation study, and an official facility update answer different questions.

The implementation boundary is equally important. This is not an experimental protocol. It contains no target design, operating setpoint, pulse timing, facility control, irradiation plan, or attempt to access a high-power laser. Its \texttt{DESIGN\_ONLY} policy avoids transforming a literature-anchored computational proposal into operational guidance. Any future execution would require authorized infrastructure, applicable safety and governance review, and qualified experts in laser-plasma interaction, high-energy-density physics, target physics, numerical methods, and diagnostics.

\section{Conclusion}
Humans can plausibly create intense laser fields with selected coherence properties analogous to sunlight: broad spectrum, stochastic phase, and random or structured polarization. The available evidence supports a stronger statement than a metaphor but a narrower statement than a fusion-energy promise. In a published, scoped numerical setting, a sunlight-like field representation can degrade resonant three-wave coupling and raise an SRS threshold relative to a monochromatic comparator \cite{ma2021}. Existing polarization-smoothing work supplies an independent rationale for incoherent addition on fast interaction timescales \cite{rothenberg2000}. At the same time, overlap-induced high-speckle effects show why no generic suppression claim is justified \cite{hao2023}.

CRISP-ICF translates that mixed evidence into a rigorous research direction. Its novelty is not a claim that incoherence is universally superior. It is a disciplined protocol that treats coherence as a multi-dimensional driver design variable, forces matched comparisons, confronts adverse coupling controls, preserves uncertainty and resource accounting, and authorizes only conditional statuses. The proposal's direct answer is therefore: intense sunlight-like drivers are scientifically plausible candidates for LPI control, but they become credible ICF-relevant options only after an evidence chain demonstrates robust, matched, and transferable benefit without unacceptable system tradeoffs. The next step is a specialist-reviewed implementation of the CRISP-ICF card set in carefully bounded simulation regimes.

\appendices
\section{Evidence Coverage and Claim Boundary}
The Survey registry contains eight sources. The NIF official webpage was inspected directly for present high-level context \cite{nif_2026}. Ma \emph{et al.} was checked on the publisher page and cross-validated through two academic-search engines \cite{ma2021}. Rothenberg, Jia, and Hao were cross-validated metadata records and are used for their narrowly stated mechanisms \cite{rothenberg2000,jia2023,hao2023}. Myatt is a discovery-stage source used only to motivate a wave/coupling fidelity requirement \cite{myatt2017}. Betti--Hurricane and Atzeni--Meyer-ter-Vehn provide background terminology and ICF system context \cite{betti2016,atzeni2004}. No source is used to claim a newly generated laser field, a performed plasma calculation, a target-scale improvement, or a new fusion result.

\section{Symbols and Operational Definitions}
\begin{tabularx}{0.96\columnwidth}{lY}
\toprule
Symbol / term & Definition in this proposal \\
\midrule
$\omega_0$, $\Delta\omega$ & Central angular-frequency convention and declared spectral support. \\
$S(\omega)$, $\phi(\omega)$ & Spectrum and phase realization used in a driver card. \\
$\tau_{\phi}$, $\tau_{\mathrm{speckle}}$ & Phase-correlation and temporal-speckle descriptors, not universal constants. \\
$\mathcal{P}_{\mathrm{Stokes}}$ & Statistical representation of time-varying polarization. \\
$Q_{\mathrm{LPI}}$ & Declared regime-specific instability-risk proxy. \\
$Q_{\mathrm{delivery}}$, $Q_{\mathrm{uniformity}}$ & Planned energy-delivery and drive-uniformity proxies. \\
$Q_{\mathrm{coupling}}$ & Planned cross-beam/coupling sensitivity proxy. \\
$\epsilon_Q^{\mathrm{claim}}$ & Conservative evidence/error ledger for a claim about quantity $Q$. \\
CRISP-ICF & Coherence-Resolved Instability Suppression Protocol for ICF. \\
\bottomrule
\end{tabularx}

\begin{thebibliography}{99}
\bibitem{nif_2026} National Ignition Facility and Photon Science, ``Expanding the fusion frontier,'' Lawrence Livermore National Laboratory. [Online]. Available: \url{https://lasers.llnl.gov/}. Accessed: Sep. 1, 2026.
\bibitem{ma2021} H. H. Ma, X. F. Li, S. M. Weng, S. H. Yew, S. Kawata, P. Gibbon, Z. M. Sheng, and J. Zhang, ``Mitigating parametric instabilities in plasmas by sunlight-like lasers,'' \emph{Matter Radiat. Extremes}, vol. 6, no. 5, 055902, 2021, doi: 10.1063/5.0054653.
\bibitem{rothenberg2000} J. E. Rothenberg, ``Polarization beam smoothing for inertial confinement fusion,'' \emph{J. Appl. Phys.}, vol. 87, no. 8, pp. 3654--3662, 2000, doi: 10.1063/1.372395.
\bibitem{jia2023} X. Jia, Q. Jia, R. Yan, and J. Zheng, ``Suppressing stimulated Raman side-scattering with vector light,'' \emph{Matter Radiat. Extremes}, vol. 8, no. 5, 2023, doi: 10.1063/5.0157811.
\bibitem{hao2023} L. Hao, J. Qiu, and W. Y. Huo, ``Generation of high intensity speckles in overlapping laser beams,'' \emph{Matter Radiat. Extremes}, vol. 8, no. 2, 2023, doi: 10.1063/5.0123585.
\bibitem{myatt2017} J. F. Myatt, R. K. Follett, J. G. Shaw, D. H. Edgell, D. H. Froula, I. V. Igumenshchev, and V. N. Goncharov, ``A wave-based model for cross-beam energy transfer in direct-drive inertial confinement fusion,'' \emph{Phys. Plasmas}, vol. 24, 2017, doi: 10.1063/1.4982059.
\bibitem{betti2016} R. Betti and O. A. Hurricane, ``Inertial-confinement fusion with lasers,'' \emph{Nature Phys.}, vol. 12, pp. 435--448, 2016, doi: 10.1038/nphys3736.
\bibitem{atzeni2004} S. Atzeni and J. Meyer-ter-Vehn, \emph{The Physics of Inertial Fusion: Beam Plasma Interaction, Hydrodynamics, Hot Dense Matter}. Oxford, U.K.: Oxford Univ. Press, 2004.
\end{thebibliography}
\end{document}
